\pdfvariable suppressoptionalinfo 512
\pdfvariable trailerid {[<\CRevisionID><\CRevisionID>]}
\documentclass[10pt]{article}
\usepackage[margin=0.78in]{geometry}
\usepackage{amsmath,amssymb,amsthm,booktabs,microtype,xurl}
\usepackage{unicode-math}
\setmainfont{Latin Modern Roman}
\setmathfont{Latin Modern Math}
\newfontfamily\zhfont[BoldFont={Droid Sans Fallback}]{Droid Sans Fallback}
\newtheorem{theorem}{Theorem}
\newtheorem{lemma}{Lemma}
\newtheorem{proposition}{Proposition}
\newcommand{\F}{\mathbf F}
\newcommand{\Z}{\mathbf Z}
\newcommand{\Q}{\mathbf Q}
\newcommand{\C}{\mathbf C}
\title{Primary Gaussian Phases and All-Degree Frobenius Dynamics\protect\\
for the Elliptic Curve $y^2=x^3-x$}
\author{HCS-C382 / HEN-O366}
\date{5 September 2026}
\begin{document}
\maketitle
\begin{abstract}
For the reductions of $y^2=x^3-x$ at every odd prime, we give a
sign-complete derivation of the Frobenius eigenvalues. At split primes, the
centralizer of an intrinsic order-four automorphism places Frobenius in
the Gaussian order, and an explicit rational point of order four fixes
the primary Gaussian unit. At inert primes a character involution gives
trace zero. The resulting unordered eigenvalue pair is determined without
a fitted phase or a point-count sign convention.
\ifnum\CRevisionRound>0
We then prove irrational normalized phase in the ordinary chamber,
exact phase order four in the supersingular chamber, and classify every
extension degree attaining a Hasse endpoint. All extension-field counts
and primitive closed-point counts follow, with a complete repetition
dictionary.
\fi
\ifnum\CRevisionRound>1
Finally we close the native graded cohomological determinant, functional
equation and critical circle, and certify 4,008 prime--degree cells through
independent exact implementations. A linear zero-count law obstructs a
fixed-fibre Riemann-target divisor. The native results retain their
source-specific meaning: strict target evaluation records weak A1 and
failed A2/A3 gates. This is a rigorous classical-theorem synthesis with
reproducible evidence, without a claim of literature priority.
\fi
\end{abstract}
\noindent\textbf{Keywords:} complex multiplication; Frobenius dynamics;
Gaussian integers; elliptic curves; primary normalization;
\ifnum\CRevisionRound>0 closed points\else rational torsion\fi

\begin{center}{\zhfont\large 中文摘要}\end{center}
{\zhfont\begin{flushleft}
本文统一处理曲线 $y^2=x^3-x$ 在所有奇素数处的约化。分裂情形中，\\
四阶自同态的中心化子迫使弗罗贝尼乌斯位于高斯整数环；显式有理\\
四阶点使群阶被八整除，从而唯一确定主高斯整数的实部符号。惯性\\
情形通过二次特征的反号对合证明迹为零，整个相位约定无需拟合。
\ifnum\CRevisionRound>0
普通约化的归一化相位与圆周无理旋转对应，超奇异约化的相位恰有\\
四阶。由此分类全部扩域次数的哈塞端点，并用轨道分解与莫比乌斯\\
反演给出每个次数的闭点计数，明确区分本原轨道与重复。
\fi
\ifnum\CRevisionRound>1
完整上同调行列式同时保留零次、一次与二次因子，给出本系统的\\
函数方程和临界圆。独立精确程序核验四千零八个素数与次数单元。\\
固定特征的零点计数仅线性增长，不能提供黎曼目标的除子对应。\\
因此路线评估仍保留弱{\normalfont A1}及失败的{\normalfont A2/A3}；论文贡献是完整\\
的经典定理综合与可复算证据，不主张文献首创或目标算子实现。
\fi
\textbf{关键词：}复乘；弗罗贝尼乌斯动力学；高斯整数；椭圆曲线；\\
主元归一化；\ifnum\CRevisionRound>0 闭点计数\else 有理挠点\fi
\end{flushleft}}

\section{Frozen object, ownership, and the first advance}
\textbf{Round 0 advance:} an all-odd-prime trace formula with a proved
Gaussian-unit convention. Let $p$ be an odd prime, let $E_p$ denote the
smooth projective curve $y^2=x^3-x$ over $\F_p$, and let
\[
 F:E_p(\overline{\F}_p)\longrightarrow E_p(\overline{\F}_p),
 \qquad F(x,y)=(x^p,y^p).
\]
This formula fixes the Frobenius convention; its cohomological pullback
has eigenvalue $p$ on degree two. We put
$N_{p,n}=\#E_p(\F_{p^n})$ and $t_p=p+1-N_{p,1}$.
The prime $2$ is excluded: the integral model has discriminant $64$.
The parameter $p$ indexes reductions of one curve and is not an
autonomous real-time coordinate.

The general closed-point/zeta dictionary was already owned locally by
flow papers 4 and 6 for a projective-line Frobenius suspension. The
earlier C41 cubic-CM package considered $y^2=x^3+1$, an inert-prime trace
law and finite extension checks. The present owner is the quartic-CM
all-prime primary phase, its ordinary/supersingular distinction, and its
complete degree consequences. These results specialize classical
elliptic-curve theory; they are not presented as newly discovered CM
theorems. We use the endomorphism and finite-field results in
Milne~\cite{Milne}, with every special sign step proved below.

\begin{theorem}[Sign-complete CM trace]
If $p\equiv3\pmod4$, then $t_p=0$. If $p\equiv1\pmod4$, there are
unique integers $a,b$ satisfying
\[
 b>0,\qquad a\text{ odd},\qquad b\text{ even},\qquad
 a^2+b^2=p,\qquad a+b\equiv1\pmod4,
\]
and $t_p=2a$. The unordered Frobenius eigenvalues are $a\pm bi$ in the
split case and $\pm i\sqrt p$ in the inert case.
\end{theorem}
\begin{proof}
Let $\chi$ be the quadratic character extended by $\chi(0)=0$. Counting
the two square roots of a nonzero square gives
\[
 N_{p,1}=p+1+\sum_{x\in\F_p}\chi(x^3-x).
\]
When $p\equiv3\pmod4$, replacing $x$ by $-x$ negates the sum, so it
vanishes as an integer. This includes the good prime $p=3$.

For $p\equiv1\pmod4$, choose $j\in\F_p$ with $j^2=-1$. The
endomorphism $I(x,y)=(-x,jy)$ satisfies $I^2=[-1]$ and commutes with
$F$. By the elliptic endomorphism classification, the endomorphism
algebra containing $\Q(I)$ is either that quadratic field or a quaternion
algebra. In either case the centralizer of $I$ is $\Q(I)$: in the latter
case write the algebra as $\Q(i)\oplus\Q(i)v$ with $vi=-iv$ and compare
the two coefficients. Thus $F$ lies in $\Q(i)$. Its integrality places
it in the maximal order $\Z[i]$, so $F=a+bi$, and degree and trace give
$a^2+b^2=p$ and $t_p=2a$.

All three nonzero $2$-torsion points are rational. Therefore $F-[1]$
kills $E[2]$ and factors through the separable isogeny $[2]$. The
endomorphism $(F-[1])/2$ is again an integral element of $\Q(i)$.
It follows that $a$ is odd and $b$ even.

Now $Q=(j,j-1)$ lies on the curve because $(j-1)^2=-2j=j^3-j$.
The tangent slope is
\[
 m=\frac{3j^2-1}{2(j-1)}=j+1,
 \qquad x(2Q)=m^2-2j=0,
 \qquad y(2Q)=mj-(j-1)=0.
\]
Hence $2Q=(0,0)$, and $Q$ has order four. The independent point
$(1,0)$ of order two gives an eight-element subgroup, so
$8\mid p+1-2a$. If $b\equiv0\pmod4$ then $p\equiv1\pmod8$ and
$a\equiv1\pmod4$; if $b\equiv2\pmod4$ then $p\equiv5\pmod8$ and
$a\equiv3\pmod4$. These are precisely $a+b\equiv1\pmod4$.
Unique factorization in $\Z[i]$ fixes $|a|,|b|$ and this congruence
fixes the sign of $a$. Taking $b>0$ chooses the upper display root.
It does not select one literal prime ideal: both conjugates satisfy
the primary congruence modulo $(1+i)^3$.
\end{proof}

\begin{center}
\begin{tabular}{rrrrl}
\toprule
$p$ & $a$ & $b$ & $t_p$ & chamber\\
\midrule
3 & --- & --- & 0 & supersingular\\
5 & $-1$ & 2 & $-2$ & ordinary\\
13 & 3 & 2 & 6 & ordinary\\
17 & 1 & 4 & 2 & ordinary\\
29 & $-5$ & 2 & $-10$ & ordinary\\
\bottomrule
\end{tabular}
\end{center}
The negative real part at $p=5$ is essential: choosing the positive
representation $1+2i$ changes the curve trace. Complex conjugation,
however, changes neither trace nor the unordered eigenvalue pair.

\ifnum\CRevisionRound>0
\section{Phase rigidity, extension counts, and primitive orbits}
\textbf{Round 1 advance:} phase periodicity is classified, then all
degrees of fixed and primitive points are determined. Let
$\alpha_p$ be the eigenvalue with positive imaginary part and let
$\theta_p\in(0,\pi)$ satisfy $\alpha_p/\sqrt p=e^{i\theta_p}$.

\begin{theorem}[Phase and Hasse endpoints]
For $p\equiv1\pmod4$ the curve is ordinary and $\theta_p/\pi$ is
irrational. For $p\equiv3\pmod4$ it is supersingular and its normalized
eigenvalues have exact order four. The Hasse bound over $\F_{p^n}$ is
strict for every $n$ in the ordinary chamber; it is saturated exactly
for even $n$ in the supersingular chamber. At $n=2m$ in that chamber,
\[
 N_{p,2m}=\bigl(p^m-(-1)^m\bigr)^2.
\]
\end{theorem}
\begin{proof}
The standard ordinary criterion is $p\nmid t_p$ for a curve over
$\F_p$~\cite{Silverman}. In the split chamber, $0<|a|<\sqrt p$ gives
$p\nmid2a$; in the inert chamber the trace is zero. If
$\alpha_p/\sqrt p$ were a root of unity in the split chamber, its square
$\alpha_p/\overline{\alpha_p}$ would be a root of unity in $\Q(i)$,
hence $1,-1,i$, or $-i$. Those cases respectively force $b=0$, $a=0$,
$a=b$, or $a=-b$, all incompatible with an odd prime norm and the parity
just established. Thus the phase is irrational. The trace over degree
$n$ is $2p^{n/2}\cos(n\theta_p)$, so equality at a Hasse endpoint
would make $n\theta_p$ an integer multiple of $\pi$, which is
impossible. In the inert chamber the eigenvalues are $\pm i\sqrt p$:
the trace is zero for odd $n$ and $2(-p)^{n/2}$ for even $n$.
Substitution yields the displayed square.
\end{proof}

\begin{theorem}[All-degree ledger]
Put $S_{p,0}=2$, $S_{p,1}=t_p$, and
$S_{p,n}=t_pS_{p,n-1}-pS_{p,n-2}$. For every integer $n\geq1$,
\[
 N_{p,n}=p^n+1-S_{p,n},\qquad
 B_{p,n}=\frac1n\sum_{d\mid n}\mu(d)N_{p,n/d}\in\Z_{\geq0},
\]
where $B_{p,n}$ counts the primitive $F$-orbits of length $n$.
\end{theorem}
\begin{proof}
The Frobenius trace theorem~\cite[IV, Theorem 9.10 and Remark 9.11]{Milne}
gives $N_{p,n}=1+p^n-\alpha_p^n-\overline{\alpha_p}^n$.
The quadratic characteristic polynomial gives the recurrence.
Every geometric point has coordinates in a finite extension of
$\F_p$, and hence lies on a finite Frobenius orbit. Such orbits are
exactly the closed points. A length-$d$ orbit contributes $d$ fixed
points to $F^n$ if and only if $d\mid n$, giving
$N_{p,n}=\sum_{d\mid n}dB_{p,d}$. Möbius inversion proves the formula,
and the orbit interpretation proves integrality and nonnegativity for
every degree without a finite cutoff.
\end{proof}

The orbit convention uses forward $F$, primitive weight one and phase
zero; $F^n$ records repetitions. The cohomological CM phase is not a
complex weight attached to each closed point. The differential of the
$p$-power morphism is zero, which must not be interpreted as the
stability of a real differentiable flow. Completeness belongs to the
closed-point theorem; numerical tests cannot replace it.
Since $F^n-[1]$ has differential $-\mathrm{id}$, it is separable and its
kernel is reduced. Thus every geometric fixed point has scheme
multiplicity one, independently of the real-flow stability question.
\fi

\ifnum\CRevisionRound>1
\section{Native determinant, continuation, and divisor boundary}
\textbf{Round 2 advance:} the complete graded determinant is proved and
its source-target boundary is tested analytically and computationally.

\begin{theorem}[Native graded determinant]
As formal series and as analytic functions for $|u|<1/p$,
\[
 Z_p(u)=\exp\left(\sum_{n\geq1}\frac{N_{p,n}}n u^n\right)
 =\prod_{d\geq1}(1-u^d)^{-B_{p,d}}
 =\frac{1-t_pu+pu^2}{(1-u)(1-pu)}.
\]
It continues rationally, satisfies $Z_p(1/(pu))=Z_p(u)$, and has two
distinct zeros on $|u|=p^{-1/2}$ and simple poles at $u=1,1/p$.
For any $\ell\ne p$ the same expression is
\[
 \prod_{r=0}^2\det(1-uF^*\mid H^r_{\mathrm{et}}
   (E_{\overline{\F}_p},\Q_\ell))^{(-1)^{r+1}}.
\]
\end{theorem}
\begin{proof}
Taking the formal logarithm of the primitive product yields
$\sum_{n\geq1}(\sum_{d\mid n}dB_{p,d})u^n/n$, which is the fixed-point
series. The Frobenius eigenvalue formula sums this logarithm to the
stated rational function. The estimate
$N_{p,n}\leq p^n+1+2p^{n/2}$ proves absolute convergence for $|u|<1/p$.
The cohomological eigenvalues are $1$ on $H^0$, $\alpha_p$ and
$\overline{\alpha_p}$ on $H^1$, and $p$ on $H^2$. Substitution proves
the reciprocal identity. The norm of the conjugate eigenvalues is
$\sqrt p$ and their imaginary parts are nonzero. Thus the two numerator
zeros are distinct and cannot cancel either pole, whose moduli differ.
\end{proof}

\begin{proposition}[HEN-O366: fixed-fibre zero-count obstruction]
Under $u=p^{-s}$ the source zeros have positive-height counting function
$N_p(T)=(\log p/\pi)T+O(1)$. Therefore the fixed-fibre numerator cannot
have the completed Riemann zeta divisor up to a zero-free entire factor.
\end{proposition}
\begin{proof}
Each equation $p^{-s}=\alpha_p^{-1}$ or
$p^{-s}=\overline{\alpha_p}^{-1}$ gives one vertical arithmetic
progression with spacing $2\pi/\log p$, lying on $\operatorname{Re}s=1/2$.
The two progressions are distinct, so their combined count is as stated.
The Riemann--von Mangoldt law has leading growth
$(T/(2\pi))\log(T/(2\pi))$~\cite{Brent}. A zero-free entire
factor preserves the divisor and cannot reconcile these counting laws.
This assertion concerns each frozen fibre, not an undefined cross-prime
product or every possible arithmetic dynamical system.
\end{proof}

\section{Exact evidence and independent checks}
The producer uses the primary Gaussian formula and an integer
second-order recurrence. Its finite grid comprises all $167$ odd primes
at most $1000$ and degrees $1$ through $24$, totaling $4,008$ cells.
An independent checker imports no producer code. It counts residue
squares directly, computes the split power traces by a binomial sum,
checks each orbit decomposition, and multiplies the primitive product
through degree $24$. Its coefficients are compared with the independent
Riemann--Roch formula $N_{p,1}(p^n-1)/(p-1)$ for $n\geq1$.
For all $13$ odd primes at most $43$, two distinct quadratic-field
algorithms count $E(\F_{p^2})$. A symbolic lane checks the matrix
determinant, reciprocal convention, torsion identities, and conjugate
factorization. These are finite regression receipts supporting the
implementations; the all-prime and all-degree claims are theorems.

Three arithmetic controls are frozen. The quadratic twist
$dy^2=x^3-x$ for a nonsquare $d$ has trace $-t_p$ by its character sum.
The simpler parent $\mathbf P^1$ has counts $p^n+1$ and no $H^1$
factor. Among odd composite labels below $100$, the five prime powers
$9,25,27,49,81$ retain prime characteristic and repetition ownership;
the twenty mixed composites cannot be field characteristics. The same
native rationality and critical-circle reasoning has no automatic
Riemann-target force for these controls. In particular the parent test
shows that source rationality alone does not establish the desired
target zeros.

The hostile suite repairs payload hashes after changing mathematical
values, conventions, route levels or scope flags; independent checking
must still reject those files. Two fresh output directories give
byte-identical evidence. Every manuscript round is built twice under a
fixed epoch, with settled-log, embedded-font, text and raster checks.
The complete receipts and command outputs accompany the paper.

\section{Strict route evaluation and limits}
The source arithmetic is intrinsic to one integral elliptic curve,
giving \texttt{A0\_STRUCTURAL\_ARITHMETIC\_RELATION}. The native
closed-point formula is complete. However, the six mandatory A1 orbit
controls and a real-flow stability comparison are not supplied, so
strict evaluator v0.2.0 records \texttt{A1\_WEAK}. The native determinant
and functional equation do not match the Riemann target; the target
gates are \texttt{A2\_FAIL} and \texttt{A3\_FAIL}. Normalized
cohomological phases give only \texttt{A4\_FORMAL\_HINT}; no compatible
Hilbert space and operator domain are constructed. The overall result
is \texttt{ROUTE\_A\_EXPLORATORY}.

The work proves no target Euler product, target local arithmetic,
bad-prime factor, root number, automorphy, target functional equation,
or target-zero match. It is not a Hilbert--P\'olya operator. Route B is
not invoked. The scope literal is
\texttt{NO\_BAD\_EULER\_OR\_ROOT\_NUMBER}.
\fi

\section{Conclusion}
Rational four-torsion fixes the primary Gaussian sign uniformly at
every split odd prime, while a character involution settles every inert
prime. This proof identifies the exact source mechanism behind the
phase convention.
\ifnum\CRevisionRound>0
The resulting distinction between irrational ordinary phase and
order-four supersingular phase classifies all extension Hasse endpoints
and closes every degree of the primitive closed-point ledger.
\fi
\ifnum\CRevisionRound>1
The native graded determinant closes without dropping a cohomological
factor. Its linear zero count also gives a precise obstruction to
transferring the fixed-fibre divisor to the Riemann target.
\fi

\begin{thebibliography}{9}
\bibitem{Milne} J. S. Milne, \emph{Elliptic Curves}, 2nd ed., World
Scientific, 2020. II \S6 (endomorphisms), III Example 3.18 (quartic CM),
IV \S9, Theorems 9.4 and 9.10, Remark 9.11 and Exercise 9.12
(finite-field counts, zeta and functional equation).
\url{https://www.jmilne.org/math/Books/EC2.pdf}.
\ifnum\CRevisionRound>0
\bibitem{Silverman} J. H. Silverman, \emph{The Arithmetic of Elliptic
Curves}, 2nd ed., Springer, 2009, V \S3 (ordinary and supersingular
reduction). \url{https://doi.org/10.1007/978-0-387-09494-6}.
\fi
\ifnum\CRevisionRound>1
\bibitem{Brent} R. P. Brent, D. J. Platt and T. S. Trudgian,
Accurate estimation of sums over zeros of the Riemann zeta-function,
\emph{Mathematics of Computation} 90 (2021), 2923--2935,
\S2, Equations (5)--(7).
\url{https://doi.org/10.1090/mcom/3652}.
\fi
\end{thebibliography}
\end{document}
